% !TEX program = xelatex
\documentclass{article}
\usepackage{geometry}\geometry{margin=23mm}
\usepackage{fontspec}\setmainfont{Times New Roman}
\usepackage[spanish]{babel}
\begin{document}
\title{La tabla de medida}\author{}\date{14 de mayo de 2026}\maketitle
\begin{abstract}Dos marcas indican la posición inicial de la tabla.\end{abstract}
\tableofcontents
\section{Las marcas}\begin{quote}«Primero se lee la marca inferior».\end{quote}
\begin{figure}[ht]\centering\fbox{\rule{0pt}{8mm}\rule{33mm}{0pt}}\caption{Posición de la tabla}\end{figure}
\section{Altura}
\begin{table}[ht]\centering\begin{tabular}{cr}Marca&Altura (mm)\\\hline 1&12\\2&18\end{tabular}\caption{Altura de las marcas}\end{table}
\begin{thebibliography}{1}\bibitem{record}\emph{Cuaderno de medidas}, 2025, \chaptername~2.\end{thebibliography}
\end{document}
